\PassOptionsToPackage{hyphens}{url}
\documentclass[]{article}

\usepackage{graphicx}
\usepackage{listings}
\usepackage{xcolor}

\lstdefinelanguage{Markdown}{
    alsoletter={\#,\*,\[,\],\(,\),\!,\`,>,---},
    morekeywords={[1]\#,\#\#,\#\#\#,\#\#\#\#,\*,\*\*,\-\,\,,\`,\`\`\`},
    morekeywords={[2]!,\[,\],\(,\)},
    sensitive=false,
    morecomment=[s]{<!--}{-->},
    morestring=[b]`,
    morestring=[b]{``},
}

\lstset{
    language=Markdown,
    basicstyle=\ttfamily\small,
    keywordstyle=[1]\color{blue}\bfseries,
    keywordstyle=[2]\color{orange}\bfseries,
    stringstyle=\color{green}\ttfamily,
    commentstyle=\color{gray!70!white}\itshape,
    breaklines=true,
    columns=fullflexible,
    frame=single,
    showstringspaces=false,
    literate={---}{{\textcolor{gray!70!white}{-}\textcolor{gray!70!white}{-}\textcolor{gray!70!white}{-}}}{3}
}

\usepackage{amsmath,amssymb}
\setcounter{secnumdepth}{-\maxdimen} % remove section numbering
\usepackage{iftex}
\ifPDFTeX
  \usepackage[T1]{fontenc}
  \usepackage[utf8]{inputenc}
  \usepackage{textcomp} % provide euro and other symbols
\else % if luatex or xetex
  \usepackage{unicode-math} % this also loads fontspec
  \defaultfontfeatures{Scale=MatchLowercase}
  \defaultfontfeatures[\rmfamily]{Ligatures=TeX,Scale=1}
\fi
\usepackage{lmodern}
\ifPDFTeX\else
  % xetex/luatex font selection
\fi
% Use upquote if available, for straight quotes in verbatim environments
\IfFileExists{upquote.sty}{\usepackage{upquote}}{}
\makeatletter
\@ifundefined{KOMAClassName}{% if non-KOMA class
  \IfFileExists{parskip.sty}{%
    \usepackage{parskip}
  }{% else
    \setlength{\parindent}{0pt}
    \setlength{\parskip}{6pt plus 2pt minus 1pt}}
}{% if KOMA class
  \KOMAoptions{parskip=half}}
\makeatother
\usepackage{listings}
\newcommand{\passthrough}[1]{#1}
% \lstset{defaultdialect=[5.3]Lua}
% \lstset{defaultdialect=[x86masm]Assembler}
\usepackage{longtable,booktabs,array}
\newcounter{none} % for unnumbered tables
\usepackage{calc} % for calculating minipage widths
% Correct order of tables after \paragraph or \subparagraph
\usepackage{etoolbox}
\makeatletter
\patchcmd\longtable{\par}{\if@noskipsec\mbox{}\fi\par}{}{}
\makeatother
% Allow footnotes in longtable head/foot
\IfFileExists{footnotehyper.sty}{\usepackage{footnotehyper}}{\usepackage{footnote}}
\makesavenoteenv{longtable}
\setlength{\emergencystretch}{3em} % prevent overfull lines
\providecommand{\tightlist}{%
  \setlength{\itemsep}{0pt}\setlength{\parskip}{0pt}}
\usepackage{bookmark}
\IfFileExists{xurl.sty}{\usepackage{xurl}}{} % add URL line breaks if available
\urlstyle{same}

\usepackage[sort&compress,numbers]{natbib}
\usepackage{hyperref}
\hypersetup{
  colorlinks=true,
  citecolor=green,
  linkcolor=blue,
  urlcolor=blue,
  pdfcreator={LaTeX via pandoc}}

\usepackage{caption}
\usepackage{minted}
\usepackage{enumitem}
\setenumerate{label=(\arabic*)}
\setmonofont{Menlo}
\setcounter{secnumdepth}{3}

\newcommand{\qq}[1]{\passthrough{\lstinline'#1'}}
\newcommand{\qe}[1]{\passthrough{\lstinline!#1!}}
\newcommand{\myline}{\begin{center}\rule{0.5\linewidth}{0.5pt}\end{center}}

% ---

\title{\textbf{Design Principles and Practical Tips for Skills and Agents in OpenCode}}
\author{Song Congwei \thanks{Corresponding author: william@bimsa.cn}\\
Beijing Institute of Mathematical Sciences and Applications (BIMSA)\\
Huairou District, Beijing}
\date{}

\begin{document}

\maketitle

\noindent\textbf{Abstract}\hspace{0.5em} \textbf{OpenCode}, as an extensible programming platform
designed for large language models(LLMs), provides two core concepts:
\textbf{Skill} (functional module) and \textbf{Agent} (task scheduler).
Building on a review of the official OpenCode documentation, this paper
systematically elaborates on the design principles, naming and file
conventions, security safeguards, and practical techniques for improving
development efficiency for skills and agents. Through several
representative cases (code review, skill/agent factory, academic
writing, etc.), it demonstrates how to achieve highly cohesive, loosely
coupled, reusable, and secure modular development in real-world
projects. This paper aims to provide OpenCode (applicable to other agent-type software) developers with a
comprehensive set of engineering practice guidelines to accelerate the
development and deployment of skills and agents.

\textbf{Keywords}: OpenCode, Skill, Agent, Design Principles,
Reusability, Security

% \myline

\section{Introduction}\label{sec:introduction}

OpenCode is a programmable collaboration platform for large language
models \cite{opencode-docs,lognroll-2026}. Its core philosophy is to
achieve ``human-machine collaboration, tasks as services'' through
Skills (specialized functional modules)
\cite{agent-skills,anthropic-skills,clawhub,skill-fish,skills-sh,jiang2026,xu2026agent} and
Agents (executors responsible for scheduling, memory, and context
management) \cite{luo2025large,li2026single}. In practice, developers often face
the following challenges:

\begin{enumerate}
\item
  How to establish unified naming and file structures to facilitate team
  collaboration;
\item
  How to maintain good maintainability and reusability of skills and
  agents while preserving functional completeness;
\item
  How to enable self-evolution and self-improvement of skills, allowing
  agents to iteratively refine their own capabilities based on feedback,
  usage patterns, and performance metrics, while maintaining stability
  and avoiding unintended behavioral drift.
\item
  How to ensure security in an open execution environment, preventing
  malicious instructions or resource leaks.
\end{enumerate}

This paper addresses the above issues by combining the official OpenCode
documentation with project experience, proposing
systematic design principles and implementation techniques, and
illustrating them through concrete examples.

% \myline

\section{Design Principles}\label{sec:design-principles}

In this section, we will mainly introduce the fundamental principles of designing skills and agents. For a more in-depth understanding, please refer to \cite{xu2026agent}.

\subsection{Responsibility
Separation}\label{sec:responsibility-separation}

The difference between a skill and an agent lies not mainly in content, but in responsibilities. Please see \autoref{tab:responsibility}.

\begin{longtable}[]{@{}
  >{\raggedright\arraybackslash}p{(\linewidth - 4\tabcolsep) * \real{0.1552}}
  >{\raggedright\arraybackslash}p{(\linewidth - 4\tabcolsep) * \real{0.3793}}
  >{\raggedright\arraybackslash}p{(\linewidth - 4\tabcolsep) * \real{0.4655}}@{}}
\caption{Responsibility of skill and agent}
\label{tab:responsibility}
\\
\toprule
\begin{minipage}[b]{\linewidth}\raggedright
Module
\end{minipage} & \begin{minipage}[b]{\linewidth}\raggedright
Primary Responsibility
\end{minipage} & \begin{minipage}[b]{\linewidth}\raggedright
Key Implementation Points
\end{minipage} \\
\midrule
\endhead
\bottomrule
\endlastfoot
Skill & Perform a single business function (e.g., file reading, network
requests, text analysis). & No direct coupling with the model \\
Agent & Responsible for scheduling, memory, and context management,
chaining multiple skills into a workflow. & Use frontmatter to declare
dependent skills and subagents (see \autoref{sec:file-conventions}) \\
\end{longtable}

The design of Skills is universal for all AI agent - type software, but the design of agents may vary. OpenCode divides agents into two categories: Primary Agents and
Subagents. The former cannot be invoked by other agents, which
preliminarily prevents circular calls between agents.

\begin{quote}
\textbf{Principle}: Skills are responsible only for business
implementation, not scheduling; Agents are responsible only for
scheduling and memory, not specific business logic. This separation of
responsibilities significantly improves code reuse rates and testing
efficiency.
\end{quote}

\subsection{Naming Conventions}\label{sec:naming-conventions}

In \autoref{tab:naming}, we propose naming the skills and the agents.

\begin{longtable}[]{@{}
  >{\raggedright\arraybackslash}p{(\linewidth - 4\tabcolsep) * \real{0.1573}}
  >{\raggedright\arraybackslash}p{(\linewidth - 4\tabcolsep) * \real{0.4157}}
  >{\raggedright\arraybackslash}p{(\linewidth - 4\tabcolsep) * \real{0.4270}}@{}}
\caption{Naming conventions of skills and agents}
\label{tab:naming}
\\
\toprule
\begin{minipage}[b]{\linewidth}\raggedright
Type
\end{minipage} & \begin{minipage}[b]{\linewidth}\raggedright
Recommended Format
\end{minipage} & \begin{minipage}[b]{\linewidth}\raggedright
Example
\end{minipage} \\
\midrule
\endhead
\bottomrule
\endlastfoot
Skill folder & lowercase + underscore/hyphen, verb phrase or agentless noun phrase & \qe{find-skills}, \qe{data-analysis} \\
Agent file & lowercase + underscore, noun or noun phrase & \qe{code\_reviewer.md}, \qe{socrates.md} \\
\end{longtable}

In practice, skill folder names still tend to use agent-bearing noun
phrases (e.g., \qe{command-creator}, \qe{skill-creator}, \qe{skill-vetter}), since skills can be regarded as
agents independent of models and tools.

\begin{quote}
\textbf{Principle}: Consistently use lowercase + underscore/hyphen,
avoid camelCase or spaces, and ensure compatibility across Linux/macOS
file systems.
\end{quote}

\subsection{File Conventions}\label{sec:file-conventions}

\begin{enumerate}
\item
  \textbf{Skill Directory Structure}

\begin{minted}{text}
skills/
└─ <skill-name>             
    ├─ SKILL.md              # Main file, includes frontmatter
    ├─ assets/               # Static resources (images, examples)
    └─ scripts/              # Scripts
\end{minted}
\item
  \textbf{Agent File Structure}

\begin{minted}{text}
agents/
└─ <agent-name>.md           # Main file (includes frontmatter)
    ├─ Frontmatter           # Header information, e.g., declared skill list
    └─ instructions          # Agent-private skills
prompts/
└─ <agent-name>.txt          # Agent system prompt
\end{minted}
\item
  \textbf{Frontmatter Example} (YAML format)

\begin{minted}{yaml}
name: developer
description: Automated assistant for full-stack development, including ...
prompt: |
  You are a full-stack development assistant. Your responsibilities include:
    - Reviewing code for style, correctness, and best practices
    - Providing expert guidance on architecture and implementation
    - Running automated tests and reporting results
  Behavior rules:
    - Be concise and structured
    - Do not modify files without explicit instructions
    - Summarize suggestions in bullet points
mode: subagent
model: openai/gpt-4o
skills:
  - code-reviewer
  - code-expert
  - code-test
\end{minted}
\end{enumerate}

\begin{quote}
\textbf{Principle}: All metadata should be placed in frontmatter,
maintaining declarative and machine-readable file properties.
\end{quote}

\subsection{Security}\label{security}

For production environments with stringent security requirements, we
summarize the following five security principles in \autoref{tab:risks} as a reference. For daily work and
learning, it is sufficient to set the searchable paths for agents, which
ensures both safety and efficiency. For an in-depth study on security, see \cite{schmotz2025agent,jiang2026sok,ling2026agent,liu2026malicious,schmotz2026}

\begin{longtable}[]{@{}
  >{\raggedright\arraybackslash}p{(\linewidth - 2\tabcolsep) * \real{0.5532}}
  >{\raggedright\arraybackslash}p{(\linewidth - 2\tabcolsep) * \real{0.4468}}@{}}
\caption{Types of system risks}
\label{tab:risks}
\\
\toprule
\begin{minipage}[b]{\linewidth}\raggedright
Risk
\end{minipage} & \begin{minipage}[b]{\linewidth}\raggedright
Mitigation Measures
\end{minipage} \\
\midrule
\endhead
\bottomrule
\endlastfoot
 Arbitrary system calls & Explicitly declare available tools in the Skill and allow only whitelisted commands; use a sandbox and restrict PATH before execution. \\
 Data leakage & Validate all external inputs against JSON Schema; encrypt sensitive information (e.g., API keys) when storing; avoid injecting secrets into prompts. \\
 Persistent attacks & Prohibit write operations to paths outside \qe{<DEFAULT_DIR>}; perform path validation and file existence checks for edit and delete operations. \\
 Prompt injection & Treat user inputs, web content, and RAG retrieval results as untrusted data; prevent external content from modifying system prompts, permissions, or tool invocation rules; require secondary confirmation for high-risk operations. \\
 Skill/tool injection & Skills should not trust tool invocation suggestions from other Skills, web pages, or files; tool parameters must be validated against a Schema. \\
 Permission & Adopt a deny-by-default policy; Skills should request only the least privilege; prohibit dynamic modification of their own permissions.
\end{longtable}

\begin{quote}
\textbf{Principle}: Always conduct a security audit at the skill's entry
layer, ensuring that every external request is strictly validated before
entering the business logic.
\end{quote}

% \myline

\section{Design Techniques}\label{design-techniques}

\subsection{Variable Pre-setting}\label{sec:variable-pre-setting}

To avoid repeatedly fetching the same information across multiple
skills, variables can be preset in the agent's memory. For example:

\begin{lstlisting}[language=markdown]
## Variables
- LANGUAGE: Python
- API_KEY: <ENCRYPTED_API_KEY>
- WORK_DIR: /path/to/workdir/

The files should be saved in <WORK_DIR> by default.
\end{lstlisting}

Subsequently, use \qe{<LANGUAGE>}, \qe{<API_KEY>}, etc., in any skill to improve maintenance efficiency. \qe{<WORK_DIR>} is typically used to set the agent's working directory, thereby restricting the scope of the agent's file read and write operations.

\subsection{Custom Commands}\label{sec:custom-commands}

We design shell-style custom commands for skills/agents with the
\qq{!} prefix. For example:

\begin{lstlisting}[language=markdown]
## User Custom Commands
!save <path>   : Save the latest conversation output or its digest (not necessarily displayed in TUI) to the specified/default path
!test          : Execute unit tests and return results
!test save     : Chained commands
!python        : Simulate REPL of Python
\end{lstlisting}

To distinguish them from actual shell commands, they can be written at
the end of the prompt. An example usage:

\begin{lstlisting}[language=markdown]
User: Generate a daily schedule !save plan.md
Agent: A daily schedule has been generated and saved to plan.md.
\end{lstlisting}

OpenCode does support custom commands(\url{https://opencode.ai/docs/commands/}), but when defined in a skill/agent
configuration file, the command is exclusive to that specific
skill/agent. Additionally, regular shell commands can be executed in the
OpenCode-TUI, such as \qe{!echo "hello"}. The
\qe{!} must be entered at the beginning. In this
case, the TUI automatically switches to a simulated shell environment.
The technique presented in this section merely facilitates operations
and simplifies prompts by simulating shell commands.

\subsection{Simple Implementation of Agent Memory}\label{sec:agent-memory}

LLMs themselves do not possess implicit recurrent memory units like
RNNs (Recurrent Neural Networks); their ``memory'' is limited to the current context window. Agent
frameworks like OpenCode are responsible for organizing conversation
history into context but do not natively include long-term memory
mechanisms. The technique for implementing such a mechanism is to set up
a memory database, allowing the agent to record conversation summaries, even the profiling at the user's personality,
store them in the database, and automatically read them on the next
startup\cite{alzubi2026,du2026survey,he2026memoryarena,li2026bmam}. Meanwhile, users can actively provide feedback to the agent, which can persistently influence the agent's behavior. See \autoref{fig:memory}. Although natural language can be used to instruct the agent to store
files, we provide custom commands for the agent's convenience. Below is a simple example.

\begin{lstlisting}[language=markdown]
## Memory Directory Structure
- Memory belonging to a specific agent
  Agent memory directory, e.g., `agent-memory/`
    +-- Conversation content (Content summary, key dialogues), e.g., `chat-content.md`
    +-- User impressions, e.g., `user-impression.md`
    \-- User feedback, `user-feedback.md`
- Shared memory `shared-memory/`

## User Custom Commands
- !init <path>    : Initialize memory file (specify path <path>, or set default path `agent-memory/`)
- !load <path>    : Load memory (recommended to auto-execute on agent startup)
- !update <path>  : Update memory: summarize recent conversations (along with previous memory) and save to file, including agent's impression of the user
- !share <path>   : Share memory with all agents, i.e., save memory to the shared memory directory `shared-memory/`

If the user changes topics, !update can be triggered automatically. When the user pays attention to a certain topic multiple times, update the impression of the user.
\end{lstlisting}

The skill snippet above provides simple long-term memory functionality. Without fine‑tuning the LLM parameters, this offers a preliminary approach to reinforcement learning and self‑evolution for the agent\cite{wang2025reinforcement,xia2026}. A more complete example can be found at \cite{agent-memory-systems,hermes-agent}.

\begin{figure}[h]
  \centering
  \includegraphics[width=0.95\textwidth]{memory.png}
  \caption{The architecture diagram of memory system}
  \label{fig:memory}
\end{figure}

\subsection{Self-Management of OpenCode}\label{sec:self-management-of-opencode}

Configure OpenCode using natural language rather than commands or manual
modifications through the TUI. For example, by prompting ``Configure
Ollama API'', OpenCode will automatically configure the Ollama API. Readers can find the skill that implements this functionality in \cite{freakwill-agents}.

\subsection{Dialogue Example (Personality Distillation)}\label{sec:dialogue-example}

Pre-populating example conversations in the agent's configuration file
enables personalized interaction. Below is a dialogue example based on
the ``personality distillation'' of the ancient Greek philosopher
Socrates:

\begin{lstlisting}[language=markdown]
## Examples
User: I think courage means not being afraid of any danger.
Agent: (In Socrates' tone) Then, if a soldier charges into enemy lines to loot, is he also unafraid of danger-does that count as courage?
User: No, that's just recklessness.
Agent: So courage is not merely "fearlessness," but must be combined with "correct judgment"? Consider further-if a person faces danger but is fearless out of ignorance, can that be called courage?
User: No, ignorance is not courage. Courage should mean knowing the danger but acting for a just cause.
Agent: Very good. Then who determines the "just cause"? You yourself, the laws of the city-state, or the gods? It seems we have arrived at another crossroads...
\end{lstlisting}

This is regarded as few-shot non-parametric learning based on LLMs
\cite{brown-2020,madotto2021few,dong-2024}. The agent can automatically reuse
established style in new conversations, enhancing user experience. This
technique is currently used for personality distillation or digital
personality creation.

% \myline

\section{Tool Overview}\label{sec:tool-overview}

OpenCode primarily includes built-in tools and MCP (Model Context
Protocol). Users can define custom tools; please refer to the online
documentation.

\subsection{Built-in Tools}\label{sec:built-in-tools}

OpenCode provides a series of built-in tools for LLMs,
enabling models to directly interact with the local codebase, execute
commands, or retrieve external information. \autoref{tab:tools} is an overview of
commonly used built-in tools (see \url{https://opencode.ai/docs/tools/} for
details).

\begin{longtable}[]{@{}
  >{\raggedright\arraybackslash}p{(\linewidth - 2\tabcolsep) * \real{0.5185}}
  >{\raggedright\arraybackslash}p{(\linewidth - 2\tabcolsep) * \real{0.4815}}@{}}
\caption{Built-in tools in OpenCode}\label{tab:tools}\\
\toprule
\begin{minipage}[b]{\linewidth}\raggedright
Tool
\end{minipage} & \begin{minipage}[b]{\linewidth}\raggedright
Description
\end{minipage} \\
\midrule
\endhead
\bottomrule
\endlastfoot
\qe{bash} & Execute arbitrary shell commands in the
project environment, e.g., \qe{git status},
\qe{npm install}. \\
\qe{read} & Read file contents; supports specifying
line ranges for large files. \\
\qe{write} & Create or overwrite files (controlled
by \qe{edit} permissions). \\
\qe{edit} & Edit files based on exact string
replacement. \\
\qe{grep} & Search file contents using regular
expressions. \\
\qe{glob} & Find file paths using glob patterns. \\
\qe{webfetch} & Fetch web content by URL, suitable
for retrieving documentation or online resources. \\
\qe{websearch} & Perform web searches using Exa AI
(requires enabling the \qe{OPENCODE\_ENABLE\_EXA}
environment variable). \\
\qe{skill} & Load and return the content of a
registered Skill (\qe{SKILL.md}). \\
\end{longtable}

\subsection{MCP}\label{sec:mcp}

MCP is a plugin mechanism in OpenCode for extending tool capabilities.
Through MCP, developers can expose any external service (e.g.,
databases, REST APIs, cloud functions, etc.) to the model via a unified
JSON-RPC interface, enabling direct invocation within dialogues. This
includes services such as code audit services, browser automation, or
enterprise knowledge bases, with permissions uniformly controlled by the
\qe{permission!} field, ensuring security and
auditability.

% \myline

\section{Skill Case Studies}\label{sec:skill-case-studies}

Below are several skill examples that have garnered significant
attention. You can search for them on platform \cite{skills-sh}.

\subsection{Code Review}\label{sec:code-review}

This is a very classic example and serves as a suitable template for readers learning about skills.

\begin{lstlisting}[language=markdown]
---
name: code-review
description: Review code changes for security, performance, and correctness. Trigger with a PR URL or diff, "review this before I merge", "is this code safe?", or when checking a change for N+1 queries, injection risks, missing edge cases, or error handling gaps.
argument-hint: "<PR URL, diff, or file path>"
---

# /code-review
> Usage: `/code-review <PR URL or file path>`

## Workflow
1. Use `bash` to call `flake8`, `pylint`, `pytest` to collect static analysis and unit test results.
2. Use `grep` / `read` to search for common risks in the code (e.g., `eval`, hardcoded credentials).
3. If external references are needed, use `webfetch` to fetch security guidelines.
4. Summarize and return a structured Markdown report (see example below).

## Example Output (Markdown)

## Code Review: <Title or File>

### Summary
Overview ...

### Key Issues
| # | File       | Line | Issue                           | Severity |
|---|------------|------|---------------------------------|----------|
| 1 | utils.py   | 42   | Use of `eval`, code injection risk | Critical |

### Suggestions
| # | File       | Line | Suggestion                                  | Category |
|---|------------|------|---------------------------------------------|----------|
| 1 | utils.py   | 42   | Replace with safe parsing library, e.g., `json.loads` | Security |

### Verdict
Approve / Request Changes / Needs Discussion
\end{lstlisting}

% \myline

\subsection{Skill Factory}\label{sec:skill-factory-skill-creator}

This section showcases an extremely popular skill example. Such skills
designed to construct other skills can be called ``meta-skills.''

\begin{lstlisting}[language=markdown]
---
name: skill-creator
description: Create new skills, modify and improve existing skills, and measure skill performance. Use when users want to create a skill from scratch, edit, or optimize an existing skill, run evals to test a skill, benchmark skill performance with variance analysis, or optimize a skill's description for better triggering accuracy.
---

# Skill Creator

At a high level, the process of creating a skill goes like this:

- Decide what you want the skill to do
- Write a draft of the skill
- Create test prompts and run them
- Evaluate results qualitatively and quantitatively
- Rewrite the skill based on feedback
- Repeat until satisfied

...

## Creating a Skill

### Capturing Intent

1. What should this skill enable Claude to do?
2. When should this skill be triggered?
3. What is the expected output format?

...

### Skill Composition

skill-name/
+-- SKILL.md (required)
\-- Bundled Resources (optional)
    +-- scripts/
    +-- references/
    \-- assets/
\end{lstlisting}

% \myline

\subsection{Academic Paper Writing}\label{sec:academic-paper-writing}

This is arguably the most popular skill. This paper itself was written
with the assistance of this skill: first, a template was manually
written(\url{https://github.com/Freakwill/opencode-agents/blob/main/template-opencode.md}),
and then prompts were used to instruct an agent equipped with this skill
to write the paper according to the template. Finally, content details
and formatting issues were handled manually.

\begin{lstlisting}[language=markdown]
---
name: research-paper-writer
description: Creates formal academic research papers following IEEE/ACM formatting standards with proper structure, citations, and scholarly writing style. Use when the user asks to write a research paper, academic paper, or conference paper on any topic.
---

This skill guides the creation of formal academic research papers that ...

## Feature Overview
- Automatically generate a paper outline (title, abstract, section structure) based on a topic.
- Invoke `latex-paper-en` to generate LaTeX documents conforming to IEEE/ACM templates.
- Use `grammar-check` for grammar and fluency checking of generated text.
- Support `!add-citation <bib>` to add citations, automatically maintaining the `.bib` file.
- `!compile paper.tex` compiles to PDF and returns a download link.
\end{lstlisting}

% \myline

\section{Agent Example}\label{agent-example}

Finally, we showcase a personal assistant agent to demonstrate agent
design. The main task in agent design revolves around orchestrating
skills and tasks. The agent's own skills should be kept as simple as
possible, primarily providing information that is highly specific and
inconvenient to share.

\begin{lstlisting}[language=markdown]
---
name: assistant
description: Personal assistant sub-agent, focused on scheduling, travel, documents, email, and other personal/work assistant tasks.
mode: subagent
model: ...
prompt: Search for files only in specific directories; answer concisely and rigorously...
permission:
  skill:
    calendar-management: allow
    news-aggregation: allow
    send-email-programmatically: allow
    research-paper-writer: allow
    self-improving-agent: allow
  task:
    '*': deny
    academic-writer: allow
    developer: allow
---

# Personal Assistant

## Directories

- ~/assistant/
- ~/owner-info/

## Special Commands

- !save: Save the conversation to the working directory
- !search: Search within the working directory
...
\end{lstlisting}

This agent reuses the permission set of multiple skills and restricts
itself to invoking only two subagents,
\qe{academic-writer} and
\qe{code-reviewer}, demonstrating fine-grained
authorization and secure design techniques.

% \myline

\section{Conclusion}\label{conclusion}

This paper has systematically reviewed the design principles, naming and
file conventions, security safeguards, and practical techniques for
improving development efficiency for skills and agents in the OpenCode
platform. Through responsibility separation, naming conventions, file
structures, and security whitelisting, a highly cohesive, loosely
coupled, and reusable modular system has been achieved. Combined with
variable pre-setting, custom commands, and memory mechanisms, the degree
of workflow automation has been significantly improved. The case studies
further demonstrate real-world implementations in typical scenarios such
as code review, skill/agent auto-generation, and academic writing,
providing reliable technical references for the large-scale promotion of
future OpenCode projects.

Future work can be pursued in the following directions: 
\begin{enumerate}
  \item Memory and Evolution: Special focus on the memory mechanisms and self-evolution
capabilities of skills and agents.
  \item Security Audit Plugin: Develop a unified security audit skill for static analysis of tools/commands,
proactively capturing potential risks.
  \item Multimodal Skills: Integrate multimodal inputs including text, images, and audio to further expand
OpenCode's applicability in AI-assisted creative scenarios.
\end{enumerate}

% \myline

\bibliographystyle{unsrtnat}
\bibliography{references}

\end{document}
